\let\negmedspace\undefined
\let\negthickspace\undefined
\documentclass[journal]{IEEEtran}
\usepackage[a5paper, margin=10mm, onecolumn]{geometry}
%\usepackage{lmodern} % Ensure lmodern is loaded for pdflatex
\usepackage{tfrupee} % Include tfrupee package

\setlength{\headheight}{1cm} % Set the height of the header box
\setlength{\headsep}{0mm}     % Set the distance between the header box and the top of the text

\usepackage{gvv-book}
\usepackage{gvv}
\usepackage{cite}
\usepackage{amsmath,amssymb,amsfonts,amsthm}
\usepackage{algorithmic}
\usepackage{graphicx}
\usepackage{textcomp}
\usepackage{xcolor}
\usepackage{txfonts}
\usepackage{listings}
\usepackage{enumitem}
\usepackage{mathtools}
\usepackage{gensymb}
\usepackage{comment}
\usepackage[breaklinks=true]{hyperref}
\usepackage{tkz-euclide} 
\usepackage{listings}
\def\inputGnumericTable{}                                 
\usepackage[latin1]{inputenc}                                
\usepackage{color}                                            
\usepackage{array}                                            
\usepackage{longtable}                                       
\usepackage{calc}                                             
\usepackage{multirow}                                         
\usepackage{hhline}                                           
\usepackage{ifthen}                                           
\usepackage{lscape}
\begin{document}

\bibliographystyle{IEEEtran}


\title{12.116}
\author{AI25BTECH11012 - GARIGE UNNATHI}
% \maketitle
% \newpage
% \bigskip
{\let\newpage\relax\maketitle}


\renewcommand{\thefigure}{\theenumi}
\renewcommand{\thetable}{\theenumi}
\setlength{\intextsep}{10pt} % Space between text and floats


\numberwithin{equation}{enumi}
\numberwithin{figure}{enumi}

\vspace{-1cm}

\textbf{Question}:\\
Let $\vec{P}$ and $\vec{Q}$ be two square matrices of same size. Consider the following statements
\begin{enumerate}
    \item $\vec{P}\vec{Q}$ = 0 then $\vec{P}$=0 or $\vec{Q}$ =0 or both
    \item $\vec{P}\vec{Q}$=$\vec{I}$ then $\vec{Q}\vec{P}$ = $\vec{I}$
    \item ($\vec{P}+\vec{Q}$)$^2$ = $\vec{P}$$^2$ + 2$\vec{P}\vec{Q}$ + $\vec{Q}$$^2$
    \item ($\vec{P}-\vec{Q}$)$^2$ = $\vec{P}$$^2$ - 2$\vec{P}\vec{Q}$ + $\vec{Q}$$^2$
\end{enumerate}
Determine which statements are True and which are False\\
\textbf{Solution:}\\

\begin{enumerate}
    \item This option is False\\
    This can be proven by a simple counter example :\\
    \begin{align}
        \vec{P} = \myvec{1&0\\0&0} \quad
        \vec{Q} = \myvec{0&0\\1&0}\\
        \vec{P}\vec{Q} = \myvec{0&0\\0&0}
    \end{align}

    \item Given $\vec{P}\vec{Q}$ = $\vec{I}$ \\
     From this we can say that $\lvert P \rvert$ and $\lvert Q \rvert$ is not zero \\
     That is inverse exists
     \begin{align}
         \vec{Q} = \vec{P}^{-1}\\
         \vec{Q}\vec{P} = \vec{P}^{-1}\vec{P} = \vec{I}
     \end{align}
     Hence this option is true 

     \item \begin{align}
         (\vec{P}+\vec{Q})^2 = (\vec{P}+\vec{Q})(\vec{P}+\vec{Q})\\
          =  \vec{P}^2 + \vec{P}\vec{Q} + \vec{Q}\vec{P} + \vec{Q}^2
     \end{align}
     This equation will not be equal to $\vec{P}^2$ + 2$\vec{P}\vec{Q}$ + $\vec{Q}^2$\\
     this is because Matrix multiplication is not commutative 
     \begin{align}
         \vec{P}\vec{Q} \neq \vec{Q}\vec{P}
     \end{align}
     this option is false 

     \item similarly option 4 is also False 
\end{enumerate}
Hence only option 2 is true
\end{document}


